% PAPER_A_METHOD_CONTRACT_CURRENT.json paper_a_method_20260802_v13 SHA-256 ca9c2788261266744b88e7bbd69d253c509a154e8291ddc0cfc3de771eb4232a
\section{统计分析与功效}

\subsection{C1 后模拟 C2-B}

只模拟以下正式候选设计：
\begin{verbatim}
D8  = 4 common anchors + 4 diverse bridges
D10 = 6 common anchors + 4 diverse bridges
D12 = 6 common anchors + 6 diverse bridges
\end{verbatim}

模拟：
\begin{itemize}
  \item 每人任务数
  \item 共同 anchor 数
  \item 多样化 bridge 数
  \item unique image 数
  \item worker--task graph
  \item \(B_u\) interval width
  \item routing activation rate
  \item fallback rate
\end{itemize}

决定 C2-B 的最终结构。

\subsection{T1 功效}

模拟：
\begin{itemize}
  \item worker/image/building 方差；
  \item structurally valid rate；
  \item delivery-adjusted quality；
  \item Manual/Semi correlation；
  \item active-time missingness；
  \item risk interaction。
\end{itemize}

\subsection{V1 功效}

模拟：
\begin{itemize}
  \item profile activation；
  \item policy divergence（仅作可识别性诊断）；
  \item capacity；
  \item refusal/timeout；
  \item dynamic \(k\)；
  \item unresolved；
  \item delivery-adjusted quality；
  \item policy\(\times\)risk interaction。
\end{itemize}

\subsection{Cross-fitted replay}

按 \texttt{image/base\_task} 划 fold。

评价 fold 不得参与：
\begin{itemize}
  \item Global selection
  \item risk feature selection
  \item weight selection
  \item support threshold
  \item fallback
  \item stopping
\end{itemize}

Replay 用于开发、消融和功效，不替代前瞻 V1。

\subsection{多重终点层级}

\subsubsection{T1 确认性}

\begin{enumerate}
  \item delivery-adjusted quality 非劣，并通过 structural-validity safety gate；
  \item owner-valid active time；
  \item risk interaction。
\end{enumerate}

T1 的质量非劣 margin、结构安全门和 active-time 差值尺度在结果可见前写入冻结的
Statistical Analysis Plan 与 analysis manifest；active time 主尺度为
\(\Delta_T=T_{\mathrm{Semi}}-T_{\mathrm{Manual}}\)，log-ratio 仅作敏感性分析。
总体 active-time coverage 以及 Manual/Semi、ordinary/stress 的差异性缺失必须同时
不超过冻结容忍度，否则 active time 降级为 descriptive/missingness sensitivity。

\subsubsection{V1 确认性}

\begin{enumerate}
  \item severe failure 非劣；
  \item unresolved＋severe failure 非劣；
  \item delivery-adjusted policy quality；
  \item 标注数量或成本。
\end{enumerate}

resolved-only quality、\(k_{used}\)、过程时间和 policy\(\times\)risk interaction 为次级
结果。V1 两个安全门和质量非劣 margin 来自冻结的 Statistical Analysis Plan 与
policy analysis manifest，按上述顺序执行，不在观察到结果后改变。

\subsubsection{RQ2}

以预测关联、效应量、方向、支持范围和预测性能为主，不要求每阶段分别达到传统显著性。

\subsubsection{V2}

可选的外部支持与安全退化审计，不作为主要因果结果。

\subsection{C1 variable-k、rolling enrollment 与学习效应}

C1 所有分母按 task--condition--estimand 的 final unique-worker support 计算，并分别报告 target、
original、authorized、late-entry、outside、deficit 与 pooled excess。W034 的 timing 主分析只使用
sentinel 后 owner-valid 的 17 条授权任务；原始缺失 timing 不插补。W034 profile 必须报告
original-only 与 original+authorized sensitivity。

若 rolling enrollment 未激活，pooled 分析与 original-cohort 分析完全相同；若激活，主分析使用
预冻结规则形成的 pooled cohort，并同时报告 original-only sensitivity、late-entry 数量、版本、
窗口和 task exposure。学习效应只进行 completion order、instruction/interface version 与 provenance
的描述性/敏感性分析，不建立结果可见后选择的复杂主模型。所有选择在 Stage 3 gate 前冻结。

\subsection{可识别跨阶段模型与双 margin}

跨 C1/C2 阶段最终模型使用 building random intercept 与 task-within-building random intercept，
并保留 stage fixed effect；不同时放入吸收 stage 的 task fixed effects。聚类同时报告
\(m^{all}=(n_1-n_2)/k\) 与 \(m^{top2}=(n_1-n_2)/(n_1+n_2)\)，并分别报告
\(k_{\mathrm{LOO,medoid}}\) 和 \(k_{\mathrm{LOO,strict}}\)。

\paragraph{Contract v5 clarification.}
No C1+C2 stage effect is interpreted without a frozen cross-stage anchor or equivalent supported structure. The final model uses building and task-within-building random intercepts with a stage fixed effect, whereas C1-only uses a task fixed effect and no stage effect. The V1 quality criterion is superiority after the two delivery hierarchy gates; it is not an outcome-selected quality margin.
